\section{One straight line, drawn by a real chain}
\label{sec:fwdline}

We know the chain draws a line.  This section fits it on the book's
data and meets the two facts that shape the whole chapter: the line is
astonishingly straight, and its residuals are not noise.

With more than two strikes the line is overdetermined, and the natural
fit is least squares: choose the slope and intercept that minimize the
sum of squared residuals across strikes.  Write the fitted intercept
as $\hat a$ and the fitted slope as $\hat b$; then, reading
\eqref{eq:fwdparity} backwards as in the hand example,
\begin{equation}
\hat D=-\hat b,
\qquad
\hat F=\frac{\hat a}{\hat D},
\qquad
\hat r=-\frac{\log\hat D}{t}.
\label{eq:fwdols}
\end{equation}
Hats, here and from now on, mark estimates --- numbers computed from
quotes, as opposed to the true $(F,D)$ the market is carrying.  The
whole estimator is four lines of arithmetic; there is nothing else
inside it.

\Cref{fig:fwdline} runs it on the running node --- SPY December 2026,
the same (underlier, expiry) pair every chapter has used --- taking
mid prices from the frozen snapshot's raw chain and keeping the
\MacFwdLineNPairs\ strikes where both the call and the put quote a
live bid.  Panel~(a) shows the observable: $\Pi(K)$ against $K$,
falling from $+\$\MacFwdLinePiMaxDollars$ deep in the left wing to
$-\$\MacFwdLinePiMinAbsDollars$ at the top of the board, a range of
\$\MacFwdLineSpanDollars.  Through
that range the points hug the fitted line so tightly that the eye
cannot separate them from it: the rms residual is
\$\MacFwdLineRmsDollars, which is \MacFwdLineStraightPct\% of the
observable's range.  The open circle marks the fitted root
$\hat F=\MacFwdLineFNaive$: the strike at which a call and a put cost
the same, which is by \eqref{eq:fwdparity} the forward estimate.
Parity, a theorem about prices that nobody enforces centrally, is
visibly enforced by the market to \MacFwdLineStraightPct\% of the
observable's own range.

Panel~(b) is the same fit with the line subtracted, and it carries the
section's real lesson.  If the residuals were quote noise they would
look like the cent-scale jitter of the hand example --- a fuzz around
zero with no shape.  Instead they \emph{bow}: a smooth arch, dollars deep at both wings,
cresting at $+\$\MacFwdLineResidMaxDollars$ near the money ---
hundreds of times the \emph{tick}, the one-cent minimum increment a
quote can move by --- with every point sitting on the arch rather
than scattered around it.  A systematic shape in residuals means a
systematic error in the model of the observation --- and the model
here was equation \eqref{eq:fwdparity}, which prices \emph{European}
options.  SPY's listed options are American: each one carries the
extra right to exercise early, that right has value, and the value is
not the same for the call and the put at a given strike.  The mids on
the line are therefore each shifted by a smooth, strike-dependent
premium the European identity knows nothing about.  Chapter~7 is
devoted to measuring and removing that premium; this chapter only
needs to know it is there.

\begin{figure}[htbp]
\centering
\paperfig{\textwidth}{fig_fwd_line.pdf}
\caption{The parity line on the running node (SPY December 2026,
\MacFwdLineNPairs\ paired strikes from the frozen chain).
(a)~$\Pi(K)=\mathsf{C}-\mathsf{P}$ against strike: a straight line to
\MacFwdLineStraightPct\% of its \$\MacFwdLineSpanDollars\ range, with
the fitted root $\hat F$ circled.  (b)~The residuals bow smoothly ---
dollars, not cents, and all of one shape: the early-exercise premium
of American quotes, the subject of Chapter~7.  The naive root and the
reference implementation's resolved forward differ by
\MacFwdLineGapBp\,bp.}
\label{fig:fwdline}
\end{figure}

The bow has a consequence the figure also shows.  Because the
contamination is asymmetric --- deep in-the-money puts carry much more
early-exercise value than calls do --- it does not average away; it
tilts the fitted line and drags the root.  The reference
implementation --- the working software whose conventions the book's
figures audit --- deals with this by de-Americanizing the quotes
before regressing (Chapter~7's machinery); its resolved forward for
this node, stored with the snapshot, is $F=\MacFwdLineFResolved$,
sitting \MacFwdLineGapBp\ basis points above the naive root
$\hat F=\MacFwdLineFNaive$.  (A \emph{basis point}, bp, is $10^{-4}$
throughout: of the forward when we compare forwards, of the rate when
we compare rates --- that one is absolute --- and later a \emph{vol}
bp, $10^{-4}$ of implied volatility.)  Both vertical lines are drawn in
panel~(b); at the scale of the strike axis they almost coincide, and
that is worth noticing too --- \MacFwdLineGapBp\,bp is
\$\MacFwdLineGapDollars\ on a spot of \$\MacFwdLineSpotDollars,
invisible in price space.
\Cref{sec:fwdskew} will show the same five dollars is anything but
invisible in volatility space.  Until then the chapter works on
European boards, where \eqref{eq:fwdparity} is exact, and returns to
American data only when the estimation theory is in hand.

For that theory we need a controlled instrument: a board whose true
$(F,D)$ we \emph{know}, so that estimator errors can be measured
rather than argued about.  The chapter's \emph{running board} is
synthetic and European: spot $S=100$, rate $r=4\%$, a $1\%$ continuous
dividend yield (dividends enter properly in \cref{sec:fwddivs}), half
a year to expiry, and \MacFwdIdentNStrikes\ strikes from 70 to 130
every $2.5$, each mid priced by the Black formula at a flat $22\%$
volatility and rounded to the cent (all constants in appendix
\ref{app:fwdprotocol}).  Its true forward is
$F=Se^{(r-q_d)t}=101.51$, where $q_d$ is the dividend yield --- take
the yield-in-the-exponent form on faith for two more sections;
\cref{sec:fwddivs} derives it.  Its true discount is
$D=e^{-rt}=0.9802$.  On this board the least-squares fit recovers the
forward to \MacFwdIdentCleanFwdErrBp\,bp and the rate to
\MacFwdIdentCleanRateErrBp\,bp --- the cent rounding is the only
noise.  The interesting question is not whether it recovers
them but \emph{how unequally well} it recovers the two, and that
question gets its own section.

We now have the estimator \eqref{eq:fwdols}, a real chain where the
line holds to \MacFwdLineStraightPct\% of its range, and a clean board with known
truth.  Next: the error bars.
